\scanpage{140}
\figure{\begin{tikzpicture}[scale=.3]\draw(0,0)--(.25,-.15)--(0,-.3);\draw(-.4,-.2)--(-.5,-.25);\end{tikzpicture}}
\begin{center}- \mspage{135} -\end{center}\margin{IX}

qui montre que $\Phi$ \uline{est canoniquement isomorphe au sous-groupe de torsion de} $H^1(I,U)$.

D'autre part, de la suite exacte (11.3.3) on déduit un isomorphisme, pour tout $n=\ell^r$ :
\[
(11.3.9)\qquad H^1(I,T_n)\simeq\operatorname{Hom}(I,T_n)\xleftarrow{\sim}\operatorname{Hom}(T,T_n)\xleftarrow{\sim}\mathbb Z/n\mathbb Z
\]
(NB on fait opérer $I$ trivialement sur $T$, donc sur $T_n=T/nT$) ; d'où, pour tout $I$-$\mathbb Z_\ell$-Module fini $M$, si
\[
(11.3.10)\qquad M'=\operatorname{Hom}(M,T\otimes D_\ell)
\]
désigne son ``dual de Cartier'', un accouplement canonique
\[
(11.3.11)\qquad H^i(I,M)\otimes H^{1-i}(I,M')\longrightarrow H^1(I,T\otimes D_\ell)\simeq D_\ell,
\]
où le dernier isomorphisme est déduit des isomorphismes (11.3.9) par passage à la limite inductive sur les puissances $n$ de $\ell$. Il est bien connu que ces accouplements sont des \uline{dualités parfaites} (cf. par exemple SGA 5 I 5). Prenant dans ces accouplements $M$ de la forme $U_n$, donc $M'=U'_n$, on trouve par passage à la limite un accouplement parfait
\[
(11.3.12)\qquad H^0(I,U\otimes D_\ell)\times H^1(I,U')\longrightarrow D_\ell,
\]
où le deuxième facteur du premier membre est un $\mathbb Z_\ell$-module de type fini. Un nouveau passage à la limite dans les accouplements (11.3.12) (pour des morphismes de transition de la forme $\ell^r.\mathrm{id}$) nous donne un deuxième accouplement parfait
\[
(11.3.13)\qquad H^0(I,U\otimes\mathbb Q_\ell)\times H^1(I,U'\otimes\mathbb Q_\ell)\longrightarrow D_\ell
\]
de vectoriels de dimension finie sur $\mathbb Q_\ell$. Désignant par $i'$, $j'$ les

\margin{\figure{\begin{tikzpicture}[scale=.35]\draw[very thick](0,0)--(0,.8);\draw[very thick](.7,0)--(.7,.7);\draw[very thick](1.3,.7)--(1.5,.7)--(1.5,0)--(1.3,0);\end{tikzpicture}}}
\figure{\begin{tikzpicture}[scale=.3]\draw(0,0)--(0,.5);\draw(.7,0)--(.7,.35);\draw(1.5,.1)--(1.5,.65);\draw(3.4,.2)--(3.5,.5);\end{tikzpicture}}
\margin{\unsure{453}}\margin{453}\margin{139}
